\PassOptionsToPackage{table}{xcolor}
\documentclass[aspectratio=169,11pt,t]{beamer}

% ============================================================
% Theme and typography
% ============================================================

\usetheme{Madrid}
\usecolortheme{default}
\usefonttheme{professionalfonts}

\setbeamertemplate{navigation symbols}{}
\setbeamertemplate{headline}{}

\setbeamertemplate{footline}{%
  \leavevmode\hbox{%
  \begin{beamercolorbox}[wd=.82\paperwidth,ht=2.5ex,dp=1.0ex,leftskip=.8em]{author in head/foot}
    CSCI 6461 --- Computer Systems Architecture
  \end{beamercolorbox}%
  \begin{beamercolorbox}[wd=.18\paperwidth,ht=2.5ex,dp=1.0ex,rightskip=.8em]{date in head/foot}
    \hfill\insertframenumber/\inserttotalframenumber
  \end{beamercolorbox}}
}

\setbeamersize{
  text margin left=0.65cm,
  text margin right=0.65cm
}

\setbeamerfont{frametitle}{size=\Large, series=\bfseries}
\setbeamerfont{block title}{size=\normalsize, series=\bfseries}
\setbeamertemplate{blocks}[rounded][shadow=false]

% Block vertical padding adjustments
\addtobeamertemplate{block begin}{\vspace{0.1em}}{\vspace{0.1em}}
\addtobeamertemplate{block end}{}{\vspace{0.2em}}

% Base list item separation
\setbeamertemplate{itemize/enumerate body begin}{\setlength{\itemsep}{4pt}}
\setbeamertemplate{itemize/enumerate subbody begin}{\setlength{\itemsep}{2pt}}

% ============================================================
% Packages
% ============================================================

\usepackage[T1]{fontenc}
\usepackage{lmodern}
\IfFileExists{sourcecodepro.sty}{\usepackage[scale=0.92]{sourcecodepro}}{\usepackage{courier}}
\usepackage{amsmath}
\usepackage{booktabs}
\usepackage{array}
\usepackage{tikz}

\usetikzlibrary{
  arrows.meta,
  positioning,
  shapes.geometric,
  calc
}

% ============================================================
% Colors
% ============================================================

\definecolor{navy}{RGB}{24,43,68}
\definecolor{deepblue}{RGB}{45,88,130}
\definecolor{lightblue}{RGB}{235,242,250}
\definecolor{softgray}{RGB}{245,247,250}
\definecolor{darkgray}{RGB}{25,28,32}
\definecolor{gold}{RGB}{200,145,35}

\setbeamercolor{palette primary}{bg=navy,fg=white}
\setbeamercolor{palette secondary}{bg=deepblue,fg=white}
\setbeamercolor{frametitle}{bg=navy,fg=white}
\setbeamercolor{title}{fg=white}
\setbeamercolor{subtitle}{fg=white}
\setbeamercolor{author}{fg=white}
\setbeamercolor{date}{fg=white}
\setbeamercolor{titlelike}{fg=white}
\setbeamercolor{block title}{bg=deepblue,fg=white}
\setbeamercolor{block body}{bg=lightblue,fg=darkgray}
\setbeamercolor{normal text}{fg=darkgray,bg=white}

% ============================================================
% Formatting shortcuts
% ============================================================

\setlength{\parskip}{0pt}

\newcommand{\key}[1]{\textbf{#1}}
\newcommand{\bits}[1]{\texttt{#1}}
\newcommand{\hex}[1]{\texttt{0x#1}}
\newcommand{\code}[1]{\texttt{#1}}

% Section divider macro
\newcommand{\sectiondivider}[2]{%

\begin{frame}[plain]

\begin{tikzpicture}[remember picture,overlay]
\fill[navy] (current page.south west) rectangle (current page.north east);
\fill[gold] (current page.south west)
  rectangle ([yshift=.15cm]current page.south east);
\end{tikzpicture}

\vspace*{\fill}

\begin{minipage}{0.94\textwidth}
\raggedright
\hyphenpenalty=10000
\exhyphenpenalty=10000

{\color{gold}\large\bfseries #1\par}

\vspace{0.3cm}

{\color{white}\fontsize{22}{26}\selectfont\bfseries
#2\par}

\end{minipage}

\vspace*{\fill}

\end{frame}
}

% ============================================================
% Metadata
% ============================================================

\title[Topic 4]{
  AArch64 Assembly Programming:\\
  Arithmetic, Logic, \& Control Flow
}

\subtitle{
  ALU mechanics, shifted-register operands, bitwise operations,\\
  condition flags, structured branches, and conditional selection
}

\author{
  Computer Systems Architecture
}

\date{}

% ============================================================
% Document
% ============================================================

\begin{document}





% ============================================================
% Slide 1
% ============================================================

\begin{frame}[plain]

\begin{tikzpicture}[remember picture,overlay]
\fill[navy] (current page.south west) rectangle (current page.north east);
\fill[gold] (current page.south west)
  rectangle ([yshift=.15cm]current page.south east);
\end{tikzpicture}

\vspace*{\fill}

\begin{minipage}{0.88\textwidth}

{\color{gold}\normalsize\bfseries
CSCI 6461 \;|\; Computer Systems Architecture \;|\; TOPIC 4
\par}

\vspace{0.25cm}

{\color{white}\fontsize{24}{28}\selectfont\bfseries
AArch64 Assembly Programming:\\[0.08cm]Arithmetic, Logic, \& Control Flow
\par}

\vspace{0.35cm}

{\color{white!80}\large
ALU mechanics, shifted-register operands, bitwise operations,\\condition flags, structured branches, and conditional selection
\par}

\end{minipage}

\vfill

{\color{white!70}\small
Dr. J. Burns
\par}

\vspace*{\fill}

\end{frame}






% ============================================================
% Slide 2
% ============================================================

\begin{frame}{Topic 4 Outline}

\setbeamertemplate{enumerate item}{\arabic{enumi}.}
\begin{enumerate}\setlength{\itemsep}{12pt}
  \item \key{ALU Operations: Arithmetic, Shifting, \& Bitwise Logic}
  \begin{itemize}
    \item Addition, subtraction, multiplication, division, and shifted-register operands.
    \item Bitwise logic with \code{and}, \code{orr}, \code{eor}, and \code{bic}.
  \end{itemize}

  \item \key{Condition Flags \& Comparisons}
  \begin{itemize}
    \item NZCV flags, \code{cmp}, and the signed vs.\ unsigned branch distinction.
  \end{itemize}

  \item \key{Control Flow: Branches \& Selection}
  \begin{itemize}
    \item Translating \code{if-else} structures and loops.
    \item Register zero tests with \code{cbz}/\code{cbnz} and conditional selection with \code{csel}.
  \end{itemize}
\end{enumerate}

\end{frame}



% ============================================================
% Section 1 Divider (Consolidated ALU: Arithmetic + Logic)
% ============================================================

\sectiondivider{PART 1}{ALU Operations: Arithmetic, Shifting, \& Bitwise Logic}





% ============================================================
% Slide 3
% ============================================================

\begin{frame}{Arithmetic Primitives: Addition \& Subtraction}
\large
\begin{itemize}\setlength{\itemsep}{6pt}
  \item Arithmetic instructions compute integer results in registers.
  \item Some arithmetic instructions can also update \key{condition flags}: processor-state bits that record properties of the result for later comparisons and branches.
\end{itemize}

\vspace{0.15cm}
\begin{center}
\small
\setlength{\tabcolsep}{6pt}
\rowcolors{2}{softgray}{white}
\begin{tabular}{lll}
\toprule
\textbf{Instruction} & \textbf{Operation} & \textbf{Sets condition flags?} \\
\midrule
\code{add  x0, x1, x2}  & $x_0 \leftarrow x_1 + x_2$ & No \\
\code{adds x0, x1, x2}  & $x_0 \leftarrow x_1 + x_2$ & \key{Yes} \\
\code{sub  x0, x1, x2}  & $x_0 \leftarrow x_1 - x_2$ & No \\
\code{subs x0, x1, x2}  & $x_0 \leftarrow x_1 - x_2$ & \key{Yes} \\
\code{neg  x0, x1}      & $x_0 \leftarrow 0 - x_1$   & No; alias for \code{sub x0, xzr, x1} \\
\bottomrule
\end{tabular}
\rowcolors{1}{}{}
\end{center}
\end{frame}





% ============================================================
% Slide 4
% ============================================================

\begin{frame}{Shifted-Register Operands and Shift Operations}
\begin{itemize}\setlength{\itemsep}{5pt}
  \item Many AArch64 arithmetic instructions can apply a shift to one register operand as part of the operation.
  \item This is useful for simple scaling, such as multiplying an index by a power of two.
\end{itemize}

\vspace{0.12cm}
\begin{center}
\scriptsize
\setlength{\tabcolsep}{4pt}
\rowcolors{2}{softgray}{white}
\begin{tabular}{p{0.22\textwidth}p{0.15\textwidth}p{0.55\textwidth}}
\toprule
\textbf{Shift Type} & \textbf{Mnemonic} & \textbf{Operation Meaning} \\
\midrule
Logical Shift Left & \code{lsl \#k} & Multiply by $2^k$ modulo the register width. \\
Logical Shift Right & \code{lsr \#k} & Unsigned division by $2^k$; zeros enter the high bits. \\
Arithmetic Shift Right & \code{asr \#k} & Signed right shift; the sign bit is copied into the high bits. \\
\bottomrule
\end{tabular}
\rowcolors{1}{}{}
\end{center}

\vspace{0.10cm}
\begin{block}{Integrated Shift-Add Examples}
\footnotesize
\rowcolors{1}{}{}
\begin{tabular}{ll}
\code{add x0, x1, x2, lsl \#3} & $x_0 \leftarrow x_1 + (x_2 \times 8)$ \\
\code{sub x0, x1, x2, lsr \#1} & $x_0 \leftarrow x_1 - \lfloor x_2/2 \rfloor$ for unsigned $x_2$ \\
\end{tabular}
\rowcolors{1}{}{}
\end{block}
\end{frame}






% ============================================================
% Slide 5
% ============================================================

\begin{frame}{Integer Multiplication \& Division}
\large
\begin{itemize}\setlength{\itemsep}{6pt}
  \item AArch64 provides separate integer instructions for multiplication and division.
  \item Division instructions return an integer quotient; any remainder is discarded.
\end{itemize}

\vspace{0.18cm}
\begin{center}
\small
\setlength{\tabcolsep}{8pt}
\rowcolors{2}{softgray}{white}
\begin{tabular}{ll}
\toprule
\textbf{Instruction} & \textbf{Purpose} \\
\midrule
\code{mul  x0, x1, x2} & Multiply two registers; keep the low 64 bits. \\
\code{sdiv x0, x1, x2} & Signed integer division. \\
\code{udiv x0, x1, x2} & Unsigned integer division. \\
\bottomrule
\end{tabular}
\rowcolors{1}{}{}
\end{center}

\vspace{0.25cm}
\normalsize
Choose \code{sdiv} or \code{udiv} based on how the operands should be interpreted: signed values use \code{sdiv}; unsigned counts, sizes, and addresses usually use \code{udiv}.
\end{frame}






% ============================================================
% Slide 6
% ============================================================

\begin{frame}{Bitwise Logic Primitives}
\large
\begin{itemize}\setlength{\itemsep}{6pt}
  \item Bitwise operations apply independently to each corresponding bit position.
\end{itemize}

\vspace{0.15cm}
\begin{center}
\small
\setlength{\tabcolsep}{6pt}
\rowcolors{2}{softgray}{white}
\begin{tabular}{lll}
\toprule
\textbf{Instruction} & \textbf{Operation} & \textbf{Bitwise Formula} \\
\midrule
\code{and x0, x1, x2} & Bitwise AND & $x_0 \leftarrow x_1 \land x_2$ \\
\code{orr x0, x1, x2} & Bitwise OR  & $x_0 \leftarrow x_1 \lor x_2$ \\
\code{eor x0, x1, x2} & Bitwise XOR & $x_0 \leftarrow x_1 \oplus x_2$ \\
\code{bic x0, x1, x2} & Bit Clear & $x_0 \leftarrow x_1 \land (\sim x_2)$ \\
\bottomrule
\end{tabular}
\rowcolors{1}{}{}
\end{center}

\vspace{0.25cm}
\normalsize
\begin{block}{Bit Clear (\texttt{bic}) Idiom}
\code{bic x0, x1, x2} clears bits in $x_1$ wherever the mask in $x_2$ has \bits{1}s. A common use is clearing a flag or disabling an option stored in one bit, without disturbing the other settings packed into the same value.
\end{block}
\end{frame}






% ============================================================
% Section 2 Divider
% ============================================================

\sectiondivider{PART 2}{Condition Flags \& Comparison Mechanics}





% ============================================================
% Slide 8
% ============================================================

\begin{frame}{NZCV Condition Flags for Control Flow}

\begin{itemize}\setlength{\itemsep}{5pt}
  \item Topic 4 is where we use the \bits{pstate} condition flags in assembly programs.
  \item The NZCV bits record the status of the most recent instruction that updates them.
  \item Common flag-setting forms include instructions with an ``\code{s}'' suffix and explicit comparisons.
\end{itemize}

\vspace{-0.2cm}
\begin{center}
\small
\setlength{\tabcolsep}{6pt}
\rowcolors{2}{softgray}{white}
\begin{tabular}{cll}
\toprule
\textbf{Flag} & \textbf{Name} & \textbf{Condition Set to 1} \\
\midrule
\key{N} & \textbf{Negative} & The most-significant bit of the result is 1. \\
\key{Z} & \textbf{Zero}     & Result is exactly zero. \\
\key{C} & \textbf{Carry}    & Unsigned carry-out (addition) or \key{no borrow} (subtraction). \\
\key{V} & \textbf{Overflow} & The signed mathematical result cannot be represented in the destination width. \\
\bottomrule
\tabularnewline
\end{tabular}
\end{center}

\vspace{-0.5cm}
\begin{block}{Carry Flag in Subtraction}
In AArch64 subtraction ($A - B$), $\text{C} = 1$ indicates $A \ge B$ (no borrow required). $\text{C} = 0$ indicates $A < B$ (unsigned borrow occurred).
\end{block}
\end{frame}





% ============================================================
% Slide 9
% ============================================================

\begin{frame}{Reading NZCV After Arithmetic}
\begin{itemize}\setlength{\itemsep}{6pt}
  \item Read the flags as a short summary of the \key{last flag-setting operation}.
  \item For \code{cmp x0, x1}, treat the flags as describing the subtraction \code{x0 - x1}; the numeric result is discarded.
\end{itemize}

\vspace{0.25cm}
\begin{block}{Simple Examples}
\begin{itemize}\setlength{\itemsep}{5pt}
  \item If \code{x0 == x1}, then \code{cmp x0, x1} produces zero, so \key{Z = 1}.
  \item If \code{x0 < x1} in an unsigned subtraction, a borrow is needed, so \key{C = 0}.
  \item If a result has its sign bit set, \key{N = 1}.
  \item If a signed result cannot fit, such as 32-bit \code{0x7fffffff + 1}, \key{V = 1}.
\end{itemize}
\end{block}
\end{frame}






% ============================================================
% Slide 10
% ============================================================

\begin{frame}{Comparison with \texttt{cmp}}
\begin{itemize}\setlength{\itemsep}{5pt}
  \item \code{cmp} compares two operands by subtracting one from the other and updating the condition flags.
  \item \code{cmp x0, x1} is really \code{subs xzr, x0, x1}: subtract, update NZCV, and discard the numerical result in the zero register.
\end{itemize}

\vspace{0.20cm}
\begin{center}
\small
\setlength{\tabcolsep}{6pt}
\rowcolors{2}{softgray}{white}
\begin{tabular}{p{0.24\textwidth}p{0.35\textwidth}p{0.32\textwidth}}
\toprule
\textbf{Instruction} & \textbf{Equivalent Operation} & \textbf{Meaning} \\
\midrule
\code{cmp x0, x1} & \code{subs xzr, x0, x1} & Compare $x_0$ with $x_1$. \\
\code{cmp x0, \#0} & \code{subs xzr, x0, \#0} & Test whether $x_0$ is zero, negative, or positive. \\
\bottomrule
\end{tabular}
\rowcolors{1}{}{}
\end{center}

\vspace{0.18cm}
After \code{cmp x0, x1}, a following conditional branch interprets the flags according to the intended signed or unsigned comparison.
\end{frame}






% ============================================================
% Slide 11
% ============================================================

\begin{frame}{Condition Codes: Signed vs.\ Unsigned Branches}
\vspace{-0.35cm}
\begin{center}
\small
\setlength{\tabcolsep}{6pt}
\rowcolors{2}{softgray}{white}
\begin{tabular}{llll}
\toprule
\textbf{Branch} & \textbf{Meaning} & \textbf{Flag Condition} & \textbf{Comparison Context} \\
\midrule
\code{b.eq} & Equal                     & $\text{Z} == 1$           & Signed / Unsigned \\
\code{b.ne} & Not Equal                 & $\text{Z} == 0$           & Signed / Unsigned \\
\midrule
\code{b.lt} & Signed Less Than ($<$)    & $\text{N} \neq \text{V}$  & \key{Signed} \\
\code{b.le} & Signed Less or Equal ($\le$) & $(\text{Z} == 1) \lor (\text{N} \neq \text{V})$ & \key{Signed} \\
\code{b.gt} & Signed Greater Than ($>$) & $(\text{Z} == 0) \land (\text{N} == \text{V})$ & \key{Signed} \\
\code{b.ge} & Signed Greater or Equal ($\ge$) & $\text{N} == \text{V}$ & \key{Signed} \\
\midrule
\code{b.lo} & Unsigned Lower ($<$)      & $\text{C} == 0$           & \key{Unsigned} \\
\code{b.hs} & Unsigned Higher or Same ($\ge$) & $\text{C} == 1$      & \key{Unsigned} \\
\code{b.hi} & Unsigned Higher ($>$)     & $(\text{C} == 1) \land (\text{Z} == 0)$ & \key{Unsigned} \\
\code{b.ls} & Unsigned Lower or Equal ($\le$) & $(\text{C} == 0) \lor (\text{Z} == 1)$ & \key{Unsigned} \\
\bottomrule
\end{tabular}
\rowcolors{1}{}{}
\end{center}

\vspace{-0.35cm}
\normalsize
\begin{block}{Usage Context}
Unsigned conditions are appropriate when the program semantics treat the values as unsigned; they are also commonly used for address and range comparisons.
\end{block}
\end{frame}


% ============================================================
% Section 3 Divider
% ============================================================

\sectiondivider{PART 3}{Control Flow: Branches \& Selection}





% ============================================================
% Slide 11
% ============================================================

\begin{frame}{Compare-and-Branch Instructions: \texttt{cbz} and \texttt{cbnz}}
\begin{itemize}\setlength{\itemsep}{4pt}
  \item \key{\code{cbz} / \code{cbnz}} test a register directly and do not alter NZCV.
  \item They combine a zero test and branch in one architectural instruction.
\end{itemize}

\vspace{0.12cm}
\begin{columns}[T,onlytextwidth]
\begin{column}{0.48\textwidth}
\begin{block}{Null Pointer Validation}
\footnotesize
\rowcolors{1}{}{}
\begin{tabular}{ll}
\code{cbz x0, .Lnull} & \text{\scriptsize // branch if null} \\
\code{ldr x1, [x0]}   & \text{\scriptsize // dereference} \\
\end{tabular}
\rowcolors{1}{}{}
\end{block}
\end{column}
\begin{column}{0.48\textwidth}
\begin{block}{Loop Counter Decrement}
\footnotesize
\rowcolors{1}{}{}
\begin{tabular}{ll}
\code{.Lloop:} & \\
\code{sub  x1, x1, \#1} & \text{\scriptsize // count--} \\
\code{cbnz x1, .Lloop} & \text{\scriptsize // repeat if nonzero} \\
\end{tabular}
\rowcolors{1}{}{}
\end{block}
\end{column}
\end{columns}
\end{frame}






% ============================================================
% Slide 12
% ============================================================

\begin{frame}[fragile]{Translating \texttt{if-then-else} Blocks}
\begin{columns}[T,onlytextwidth]
\begin{column}{0.44\textwidth}
\begin{block}{High-Level C Code}
\small
\begin{verbatim}
if (x > y) {
    result = x - y;
} else {
    result = y - x;
}
\end{verbatim}
\end{block}
\end{column}

\begin{column}{0.52\textwidth}
\begin{block}{AArch64 Assembly Mapping}
\small
\rowcolors{1}{}{}
\begin{tabular}{ll}
\multicolumn{2}{l}{\textit{// Assume: x, y are signed int64\_t}} \\
\multicolumn{2}{l}{\textit{// x0=x, x1=y, x2=result}} \\
\code{cmp  x0, x1}        & \text{\scriptsize // Compare x and y} \\
\code{b.le .Lelse}        & \text{\scriptsize // Invert test: jump if x <= y} \\
\code{sub  x2, x0, x1}    & \text{\scriptsize // Then-body: result = x - y} \\
\code{b    .Lend}         & \text{\scriptsize // Skip else block} \\
\code{.Lelse:}            & \\
\code{sub  x2, x1, x0}    & \text{\scriptsize // Else-body: result = y - x} \\
\code{.Lend:}             & \\
\end{tabular}
\rowcolors{1}{}{}
\end{block}
\end{column}
\end{columns}

\vspace{0.3cm}
\normalsize
\key{Assembly Pattern:} Invert the high-level condition to fall through into the ``then'' path.
\end{frame}





% ============================================================
% Slide 13
% ============================================================

\begin{frame}[fragile]{Loop Translation: Inverted Loop Pattern}
\begin{columns}[T,totalwidth=\textwidth]
\begin{column}{0.37\textwidth}
\begin{block}{High-Level C Loop}
\small
\begin{verbatim}
// Compute array sum
int64_t i = 0;
int64_t n = ...;
while (i < n) {
    sum += arr[i];
    i++;
}
\end{verbatim}
\end{block}
\end{column}
\hfill
\begin{column}{0.58\textwidth}
\begin{block}{AArch64 Loop with Test at the End}
\scriptsize
\setlength{\tabcolsep}{2pt}
\rowcolors{1}{}{}
\begin{tabular}{ll}
\multicolumn{2}{l}{\textit{// x0=arr, x1=n, x2=i, x3=sum}} \\
\code{mov  x2, xzr}       & \text{// i = 0} \\
\code{mov  x3, xzr}       & \text{// sum = 0} \\
\code{b    .Ltest}        & \text{// test first} \\
\code{.Lbody:}            & \\
\code{ldr  x4, [x0, x2, lsl \#3]} & \text{// load arr[i]} \\
\code{add  x3, x3, x4}    & \text{// sum += arr[i]} \\
\code{add  x2, x2, \#1}   & \text{// i++} \\
\code{.Ltest:}            & \\
\code{cmp  x2, x1}        & \text{// i < n?} \\
\code{b.lt .Lbody}        & \text{// repeat} \\
\end{tabular}
\rowcolors{1}{}{}
\end{block}
\end{column}
\end{columns}

\vspace{0.25cm}
\normalsize
\key{Common Loop Layout:} After the first jump to the test, each repeat runs the body and then checks whether to loop again.
\end{frame}





% ============================================================
% Slide 14
% ============================================================

\begin{frame}{Conditional Selection with \texttt{csel}}
\small
\begin{itemize}\setlength{\itemsep}{4pt}
  \item \code{csel} chooses between two register operands using the current condition flags.
  \item It is useful for short value selections where both candidate values are already available in registers.
\end{itemize}

\vspace{0.12cm}
\begin{block}{Signed Maximum: $\code{result} = (x_1 > x_2) \mathbin{?} x_1 : x_2$}
\footnotesize
\rowcolors{1}{}{}
\begin{tabular}{ll}
\code{cmp  x1, x2}          & \text{\scriptsize // set flags from x1 - x2} \\
\code{csel x0, x1, x2, gt}  & \text{\scriptsize // if x1 > x2, choose x1; else choose x2} \\
\end{tabular}
\rowcolors{1}{}{}
\end{block}

\vspace{0.16cm}
\normalsize
The condition code on \code{csel} is interpreted the same way as a conditional branch condition.
\end{frame}


% ============================================================
% Section 4 Divider
% ============================================================

\sectiondivider{SUMMARY \& ASSESSMENT}{Review and Self-Test Questions}





% ============================================================
% Slide 15
% ============================================================

\begin{frame}{Topic 4: Summary}
\small
\begin{itemize}\setlength{\itemsep}{4pt}
  \item \key{Arithmetic \& Shifts:} AArch64 supports register arithmetic, flag-setting variants, and shifted-register operands.
  \item \key{Bitwise Operations:} \code{and}, \code{orr}, \code{eor}, and \code{bic} operate directly on packed binary data.
  \item \key{NZCV Flags:} Flag-setting instructions record result properties used by later comparisons and branches.
  \item \key{Comparisons:} \code{cmp} subtracts, updates NZCV, and discards the numerical result.
  \item \key{Control Flow:} Conditional branches, \code{cbz}/\code{cbnz}, and \code{csel} translate common decisions and loops.
\end{itemize}
\end{frame}






% ============================================================
% Slide 16
% ============================================================

\begin{frame}{Self-Test Questions: Arithmetic \& Logic}

\begin{enumerate}\setlength{\itemsep}{7pt}
  \item \key{Shifted Addition:} What is the exact mathematical operation and resulting decimal value stored in \bits{x0} after executing:
  $$\code{mov x1, \#12} \quad \code{mov x2, \#5} \quad \code{add x0, x1, x2, lsl \#2}$$

  \item \key{Bit Clear Operation:} If register \bits{x1} contains \hex{ff} and \bits{x2} contains \hex{0f}, what value resides in \bits{x0} after executing \code{bic x0, x1, x2}?

  \item \key{Signed vs.\ Unsigned Comparisons:} If \bits{x0} contains $-1$ (\hex{ffffffff\_ffffffff}) and \bits{x1} contains $1$, which branch will be taken after \code{cmp x0, x1}: \code{b.lt} or \code{b.lo}? Why?

  \item \key{Flag Updating:} Explain why \code{add x0, x1, x2} does not update the Zero ($Z$) flag even if the resulting sum is exactly zero.
\end{enumerate}

\end{frame}





% ============================================================
% Slide 17
% ============================================================

\begin{frame}{Self-Test Questions: Control Flow \& Selection}

\begin{enumerate}\setlength{\itemsep}{7pt}
  \item \key{Loop Translation:} In the loop example, why does the code jump to the test before the first iteration?

  \item \key{Branchless Translation:} Write an equivalent 3-instruction branchless AArch64 sequence using \code{cmp}, \code{neg}, and \code{csel} to implement:
  $$\code{int64\_t abs\_val = (val < 0) ? -val : val;}$$

  \item \key{Compare-and-Branch Instructions:} When should an assembly programmer use \code{cbz x0, label} instead of \code{cmp x0, \#0} followed by \code{b.eq label}?

  \item \key{Signed vs.\ Unsigned Branches:} Why are \code{b.lt} and \code{b.lo} not interchangeable after the same \code{cmp} instruction?
\end{enumerate}

\end{frame}



% ============================================================
% Look Ahead to Topic 5
% ============================================================






% ============================================================
% Slide 18
% ============================================================

\begin{frame}{Looking Ahead}
\vfill
\begin{block}{Next Lecture Preview}
\key{Our Next Topic:} Memory Access \& AArch64 Addressing Modes
\end{block}

\vfill
\begin{itemize}\setlength{\itemsep}{7pt}
  \item \key{Load/Store Discipline:} Moving values between registers and memory.
  \item \key{Transfer Widths:} Working with byte, halfword, word, and doubleword data.
  \item \key{Effective Addresses:} Forming addresses from base registers, offsets, and indexes.
  \item \key{Arrays \& Pointers:} Applying addressing modes to common memory-access patterns.
\end{itemize}
\vfill
\end{frame}


\end{document}